\addcontentsline{toc}{subsubsection}{Solution 65. Three Hidden Primes}

\begin{dlsummary}[title={\bfseries Solution 65. Three Hidden Primes},breakable]

\begin{proof}
For a prime \(p\), let \(v_p(x)\) denote the exponent of \(p\) in a nonzero rational \(x\), allowing negative exponents.

The three solutions are
\[
\boxed{
f(0)=0,\qquad f(x)=30^{-v_p(x)}\quad(x\ne0),
\qquad p\in\{2,3,5\}.
}
\]

\paragraph{1. Establish the basic identities.}

Since \(f(30)>0\), multiplicativity forces \(f(1)=1\). For \(x\ne0\),
\[
f(x)f(1/x)=1.
\]
Thus nonzero inputs have positive values.

Also,
\[
f(-1)^2=f(1)=1,
\]
so \(f(-1)=1\) and
\[
f(-x)=f(x). \tag{1}
\]

\paragraph{2. Obtain a uniform bound on integers.}

Let
\[
M=\max\{f(0),f(1),\ldots,f(29)\}.
\]
Write a positive integer \(N\) in base \(30\):
\[
N=d_0+d_1 30+\cdots+d_s30^s,
\qquad 0\le d_j\le29.
\]
Subadditivity and multiplicativity give
\[
f(N)
\le\sum_{j=0}^{s}f(d_j)f(30)^j
\le M\sum_{j=0}^{s}30^{-j}
<\frac{30M}{29}.
\]
Hence \(f\) is uniformly bounded on the positive integers.

Now apply this bound to \(N=n^r\):
\[
f(n)^r=f(n^r)<\frac{30M}{29}.
\]
Since this holds for every positive integer \(r\),
\[
\boxed{f(n)\le1\qquad(n\in\mathbb Z_{>0}).} \tag{2}
\]
By (1), the same bound holds for negative integers.

\paragraph{3. Locate a prime with value below \(1\).}

Multiplicativity gives
\[
f(2)f(3)f(5)=f(30)=\frac1{30}.
\]
All three factors belong to \((0,1]\), so at least one is below \(1\).

We next prove that \textbf{at most one prime anywhere} can have value below \(1\).

\paragraph{4. Rule out two exceptional primes.}

Suppose distinct primes \(p,q\) satisfy
\[
f(p)<1,\qquad f(q)<1.
\]
Choose a positive integer \(r\) sufficiently large that
\[
f(p)^r+f(q)^r<1.
\]

Since \(p^r\) and \(q^r\) are coprime, Bézout’s identity supplies integers \(a,b\) such that
\[
ap^r+bq^r=1.
\]
Using subadditivity, multiplicativity, and the integer bound,
\[
\begin{aligned}
1=f(1)
&\le f(a)f(p)^r+f(b)f(q)^r\\
&\le f(p)^r+f(q)^r\\
&<1,
\end{aligned}
\]
a contradiction.

Therefore exactly one prime \(p\) has \(f(p)<1\). It belongs to \(\{2,3,5\}\); every other prime has value \(1\). Consequently,
\[
f(p)=\frac1{30}.
\]

\paragraph{5. Determine every rational value.}

Unique prime factorisation and multiplicativity now give
\[
f(x)=30^{-v_p(x)}
\]
for every nonzero rational \(x\).

\paragraph{6. Verify the three candidates.}

Multiplicativity and the normalisations are immediate.

For nonzero \(x,y\) with \(x+y\ne0\),
\[
v_p(x+y)\ge\min\{v_p(x),v_p(y)\}.
\]
Therefore
\[
f(x+y)
\le\max\{f(x),f(y)\}
\le f(x)+f(y).
\]
Cases involving zero or \(x+y=0\) are immediate. All three functions work.
\end{proof}
\end{dlsummary}

\begin{takeaways}
\begin{itemize}
\item \textbf{A shrinking base controls all integers:} Base-\(30\) expansions turn \(f(30)<1\) into a uniform bound.
\item \textbf{Powers sharpen boundedness:} Uniformly bounded values of \(f(n)^r\) force \(f(n)\le1\).
\item \textbf{Bézout connects unrelated primes:} Two independently shrinking primes would make the value at \(1\) too small.
\item \textbf{One normalisation can leave several solutions:} The factorisation \(30=2\cdot3\cdot5\) creates exactly three possible controlling primes.
\end{itemize}
\end{takeaways}
